% ============================================================
%  Chapter 6 — Matrices
% ============================================================
\section{Matrices}

Matrices encode linear maps between finite-dimensional vector spaces and are
the primary computational tool of linear algebra.  This chapter covers the
arithmetic of matrices, row reduction, systems of linear equations, and matrix
invertibility.

% ---- Basic definitions -------------------------------------
\begin{defbox}[Definition --- Matrix over $\F$]
An $m\times n$ \textbf{matrix} $A$ over $\F$ is a rectangular array of $mn$
entries $a_{ij}\in\F$, with $m$ rows and $n$ columns.  The $(i,j)$-entry is
$a_{ij}$.

\begin{itemize}
    \item The set of all $m\times n$ matrices over $\F$ is denoted
          $\Mmn{m}{n}$ (also written $\F^{m\times n}$).
    \item A \emph{square} matrix has $m=n$; its set is $\Mn{n}$.
    \item \emph{Column representation:}
          $A = [u_1 \mid u_2 \mid \cdots \mid u_n]$
          where $u_j\in\F^m$ is the $j$-th column of $A$.
\end{itemize}
\end{defbox}

% ---- Arithmetic --------------------------------------------
\begin{defbox}[Definition --- Matrix Arithmetic]
For $A,B\in\Mmn{m}{n}$, $c\in\F$, and $\mathbf{v}\in\F^n$:
\begin{itemize}
    \item \emph{Addition:} $(A+B)_{ij} = a_{ij}+b_{ij}$.
    \item \emph{Scalar multiplication:} $(cA)_{ij} = ca_{ij}$.
    \item \emph{Matrix–vector product:}
          $(A\mathbf{v})_i = \displaystyle\sum_{j=1}^n a_{ij}v_j$.
          Equivalently, $A\mathbf{v} = v_1(Ae_1)+\cdots+v_n(Ae_n)$.
\end{itemize}
\end{defbox}

% ------------------------------------------------------------
\begin{propbox}[Proposition --- Properties of Matrix–Vector Multiplication]
For $A\in\Mmn{m}{n}$, $\mathbf{v},\mathbf{w}\in\F^n$, $c\in\F$:
\[
    A(\mathbf{v}+\mathbf{w}) = A\mathbf{v}+A\mathbf{w}, \qquad
    A(c\mathbf{v}) = c(A\mathbf{v}).
\]
\end{propbox}

% ---- Matrix multiplication ---------------------------------
\begin{defbox}[Definition --- Matrix Multiplication]
For $A\in\Mmn{m}{n}$ and $B\in\Mmn{n}{p}$, the \textbf{product} $C=AB\in\Mmn{m}{p}$
has entries
\[
    c_{ij} = \sum_{r=1}^{n} a_{ir}\,b_{rj}.
\]
Equivalently, $AB = [A(Be_1)\mid A(Be_2)\mid\cdots\mid A(Be_p)]$.
\end{defbox}

% ------------------------------------------------------------
\begin{exbox}[Example --- Matrix Multiplication]
Let
\[
    A = \begin{pmatrix}1&2&3\\-2&-1&-1\end{pmatrix}, \qquad
    B = \begin{pmatrix}0&1&-1&1\\-1&0&1&-1\\1&0&-1&1\end{pmatrix}.
\]
Then
\[
    AB = \begin{pmatrix}1&1&-2&2\\0&-2&2&-2\end{pmatrix}.
\]
\end{exbox}

% ------------------------------------------------------------
\begin{propbox}[Proposition --- Properties of Matrix Multiplication]
\begin{enumerate}
    \item \emph{Associativity:} $A(BC)=(AB)C$.
    \item \emph{Distributivity:} $A(B+C)=AB+AC$; $(A+B)C=AC+BC$.
    \item \emph{Identity:} $I_m A = A = A I_n$ for $A\in\Mmn{m}{n}$.
    \item \emph{Zero:} $O_{m\times n}A = O_{m\times p}$ and
          $AO_{n\times p}=O_{m\times p}$.
    \item \emph{Powers:} $A^0=I_n$ and $A^k=A\cdot A^{k-1}$ for $A\in\Mn{n}$.
    \item \emph{Transpose:} $(AB)^T = B^T A^T$.
    \item \emph{Not commutative in general:} $AB\neq BA$ in general.
\end{enumerate}
\end{propbox}

% ---- Elementary row operations -----------------------------
\begin{defbox}[Definition --- Elementary Row Operations]
The three \textbf{elementary row operations} on a matrix are:
\begin{enumerate}
    \item \textbf{Type 1 — Row scaling:} $r_i \leftarrow \alpha r_i$,
          $\alpha\in\F\setminus\{0\}$.
    \item \textbf{Type 2 — Row addition:}
          $r_i \leftarrow r_i + \alpha r_j$ ($i\neq j$), $\alpha\in\F$.
    \item \textbf{Type 3 — Row swap:} $r_i \leftrightarrow r_j$.
\end{enumerate}
Two matrices are \textbf{row-equivalent} if one can be obtained from the other
by a finite sequence of elementary row operations.
\end{defbox}

% ------------------------------------------------------------
\begin{thmbox}[Theorem --- Invertibility of Elementary Row Operations]
The inverse of every elementary row operation exists and is an elementary row
operation of the same type.
\end{thmbox}

% ---- Echelon forms -----------------------------------------
\begin{defbox}[Definition --- Row-Echelon and Reduced Row-Echelon Form]
A matrix is in \textbf{row-echelon form (REF)} if:
\begin{itemize}
    \item All zero rows are at the bottom.
    \item The leading nonzero entry (pivot) of each nonzero row lies strictly
          to the right of the pivot of the row above.
\end{itemize}
It is in \textbf{reduced row-echelon form (RREF)} if, additionally:
\begin{itemize}
    \item Every pivot equals $1$.
    \item Every pivot is the only nonzero entry in its column.
\end{itemize}
\end{defbox}

% ------------------------------------------------------------
\begin{thmbox}[Theorem --- Uniqueness of RREF]
Every matrix is row-equivalent to a \emph{unique} matrix in RREF.
\end{thmbox}

% ---- Linear systems ----------------------------------------
\begin{defbox}[Definition --- Linear System and Augmented Matrix]
A \textbf{system of linear equations} $A\mathbf{x}=\mathbf{b}$ (with
$A\in\Mmn{m}{n}$, $\mathbf{b}\in\F^m$) is represented by the
\textbf{augmented matrix} $[A\mid\mathbf{b}]$.  The system is
\emph{homogeneous} if $\mathbf{b}=\mathbf{0}$.
\end{defbox}

% ------------------------------------------------------------
\begin{methodbox}[Method --- Solving a Linear System via RREF]
Let $[R\mid\mathbf{b}']$ be the RREF of $[A\mid\mathbf{b}]$.
\begin{enumerate}
    \item \textbf{Inconsistent:} if some row $r_i$ of $R$ is all zeros but
          $b'_i\neq 0$, the system has no solution.
    \item \textbf{Unique solution:} if consistent and every column of $R$ is a
          pivot column, the solution is unique.
    \item \textbf{Infinitely many solutions:} if consistent and some columns of
          $R$ are non-pivot (free) columns, those columns correspond to free
          (independent) variables; pivot columns give the dependent variables.
\end{enumerate}
\end{methodbox}

% ---- Invertible matrices -----------------------------------
\begin{defbox}[Definition --- Left, Right, and Two-Sided Inverses]
For $A\in\Mn{n}$:
\begin{itemize}
    \item $B$ is a \emph{left inverse} of $A$ if $BA=I_n$.
    \item $C$ is a \emph{right inverse} of $A$ if $AC=I_n$.
    \item $A$ is \textbf{invertible} if there exists $A^{-1}$ s.t.\
          $AA^{-1}=A^{-1}A=I_n$.
\end{itemize}
\end{defbox}

% ------------------------------------------------------------
\begin{thmbox}[Theorem --- Left and Right Inverses Coincide]
If $A$ has a left inverse $B$ and a right inverse $C$, then $B=C$.
Consequently the inverse $A^{-1}$ is unique, and $(A^{-1})^{-1}=A$.
\end{thmbox}

\begin{proofbox}[Proof --- Left and Right Inverses Coincide]
$B = B(AC) = (BA)C = C$. \hfill$\blacksquare$
\end{proofbox}

% ------------------------------------------------------------
\begin{thmbox}[Theorem --- Invertibility and Rank]
For $A\in\Mn{n}$, the following are equivalent:
\begin{enumerate}
    \item $A$ is invertible.
    \item $\rk(A) = n$.
    \item The linear map $T_A:\F^n\to\F^n$, $\mathbf{v}\mapsto A\mathbf{v}$,
          is an isomorphism.
\end{enumerate}
\end{thmbox}

% ------------------------------------------------------------
\begin{methodbox}[Method --- Computing the Inverse via Row Reduction]
To find $A^{-1}$ for $A\in\Mn{n}$: form the augmented matrix $[A\mid I_n]$
and row-reduce.  If $A$ is invertible, the result is $[I_n\mid A^{-1}]$.
\end{methodbox}

% ---- Diagonal matrices -------------------------------------
\begin{defbox}[Definition --- Diagonal Matrix]
A matrix $A\in\Mn{n}$ is \textbf{diagonal} if $a_{ij}=0$ whenever $i\neq j$.
Notation: $A = \operatorname{diag}(d_1,\ldots,d_n)$.
\end{defbox}

% ------------------------------------------------------------
\begin{propbox}[Proposition --- Invertibility of Diagonal Matrices]
$\operatorname{diag}(d_1,\ldots,d_n)$ is invertible if and only if
$d_i\neq 0$ for all $i$.  In that case,
\[
    \operatorname{diag}(d_1,\ldots,d_n)^{-1}
    = \operatorname{diag}\!\left(\tfrac{1}{d_1},\ldots,\tfrac{1}{d_n}\right).
\]
\end{propbox}
